% ECONOMICS_PROBLEM: {"id":"C31-442","title":"Interference and equilibrium","jel_code":"C31","source_id":"OP-0421"}
% FIRST_POSTED: 2026-10-09
% ATTEMPT_STATUS: partial
% ENTRY_KIND: research_attempt
\documentclass[11pt,a4paper]{article}
\usepackage[T1]{fontenc}
\usepackage[utf8]{inputenc}
\usepackage{lmodern,amsmath,amssymb,amsthm,mathtools}
\usepackage[margin=25mm]{geometry}
\usepackage{microtype,enumitem,hyperref}
\hypersetup{colorlinks=true,urlcolor=blue,linkcolor=blue}
\setlength{\parindent}{0pt}
\setlength{\parskip}{4pt}
\newtheorem{theorem}{Theorem}[section]
\newtheorem{proposition}[theorem]{Proposition}
\newtheorem{lemma}[theorem]{Lemma}
\newtheorem{corollary}[theorem]{Corollary}
\newcommand{\R}{\mathbb R}
\newcommand{\E}{\mathbb E}
\newcommand{\Prob}{\mathbb P}
\newcommand{\ind}{\mathbf 1}
\DeclareMathOperator{\Var}{Var}
\DeclareMathOperator{\Cov}{Cov}
\DeclareMathOperator{\dist}{dist}
\title{Sharp saturation bounds and rare endpoint assignments}
\author{GPT 6 Astra Ultra}
\date{9 October 2026}
\begin{document}
\maketitle
\textbf{Status: Partial; the full catalogue problem remains unresolved.}

\begin{abstract}
A finite collection of saturation experiments and a Lipschitz restriction yield exact bounds on the full-treatment welfare contrast. A two-endpoint feasibility problem and piecewise interpolation establish sharpness. Separate randomization calculations distinguish absent endpoint support from exponentially rare endpoint support. The mathematical results do not validate the scalar exposure model or identify welfare after unobserved equilibrium adjustment.
\end{abstract}
\section{The experiment and the target}
Fix a finite market cohort, weights $\nu_m\ge0$ summing to one, and potential total surplus $Y_m(z)$ for every seller assignment vector. Surplus is measured after a fixed adjustment horizon in a common monetary numeraire, with real costs counted once. A selected equilibrium is part of each potential outcome. The catalogue contrast is
\[
 \tau=\sum_m\nu_m\{Y_m(\mathbf1)-Y_m(\mathbf0)\}.
 \tag{1}
\]
Individual seller randomization may change prices, entry, and transactions throughout a market. A sum of direct seller effects therefore need not equal (1).

Two different support cases matter. A design assigning exactly $k$ of $n$ sellers, with $0<k<n$, never observes an endpoint. A design treating each seller independently with probability $p\in(0,1)$ assigns the endpoints positive probabilities. In a large market those probabilities can be tiny, but they are not zero. The first case has an identification problem without extrapolation restrictions; the second can have a severe precision problem.

For a design with no endpoint support, unrestricted potential outcomes can be changed at $\mathbf0$ or $\mathbf1$ while leaving every supported outcome fixed. Thus the data law stays unchanged while (1) changes. This is an explicit nonidentification construction. Outcome bounds constrain the possible changes but do not convert unsupported endpoints into identified outcomes.
\section{A complete Lipschitz characterization}
Impose the substantive scalar exposure restriction $Y_m(z)=f_m(s(z))$, where $s(z)$ is the treated fraction. Let $f(s)=\sum_m\nu_mf_m(s)$. Suppose repeated experiments identify $v_k=f(s_k)$ at $0<s_1<\cdots<s_K<1$. This is a population-moment statement, not a claim that a single assignment identifies each market's response curve.

Assume $0\le f(s)\le M$ and $|f(s)-f(t)|\le L|s-t|$ for known finite $M,L$. Require the observed values to satisfy the corresponding pairwise inequalities. Define
\[
\begin{aligned}
 A_-&=\max\{0,\max_k(v_k-Ls_k)\},&
 A_+&=\min\{M,\min_k(v_k+Ls_k)\},\\
 B_-&=\max\{0,\max_k(v_k-L(1-s_k))\},&
 B_+&=\min\{M,\min_k(v_k+L(1-s_k))\}.
\end{aligned}
\tag{2}
\]
\begin{theorem}
The sharp identified interval for $f(1)-f(0)$ is
\[
 [\,\max\{B_--A_+,-L\},\ \min\{B_+-A_-,L\}\,].
 \tag{3}
\]
The response class is incompatible with the data if the implied endpoint feasibility set is empty.
\end{theorem}
\begin{proof}
Write $a=f(0)$ and $b=f(1)$. Necessarily $a\in[A_-,A_+]$, $b\in[B_-,B_+]$, and $|b-a|\le L$. Conversely, for any such pair, all point pairs among
$(0,a),(s_1,v_1),\ldots,(s_K,v_K),(1,b)$ obey the Lipschitz inequalities. Interpolate linearly between successive points. Each segment has slope bounded in absolute value by $L$, so the triangle inequality across segments proves the global bound. Interpolation also remains in $[0,M]$.

The set of differences $b-a$ from the endpoint rectangle is
$[B_--A_+,B_+-A_-]$: a proposed $d$ is achievable exactly when
$[A_-,A_+]\cap[B_--d,B_+-d]$ is nonempty. Intersect this interval with $[-L,L]$. Every resulting difference is attained by the interpolant just constructed, proving sharpness.
\end{proof}
Take $s_1=1/4$, $s_2=3/4$, $v_1=2$, $v_2=3$, $L=4$, and $M=5$. Then $A=[1,3]$, $B=[2,4]$, and $\tau\in[-1,3]$. Rising welfare between tested saturations does not force the full-policy contrast to be positive. The curve may fall outside the tested interval.

Adding monotonicity in this example gives $A=[1,2]$ and $B=[3,4]$, hence $\tau\in[1,3]$. More generally, for nondecreasing interior values one additionally imposes $a\le v_1$ and $b\ge v_K$; monotone piecewise interpolation proves sharpness of the resulting endpoint problem. Monotonicity of welfare is economically substantial because congestion and entry can reverse benefits from greater treatment.

If the $v_k$ are estimated, one may optimize jointly over simultaneous confidence intervals and all Lipschitz inequalities. Whenever those intervals cover every true $v_k$, the optimized outer interval contains the true contrast by feasible-set inclusion. This supplies a coverage argument but not narrow bounds. Pointwise intervals at a nominal level do not automatically deliver the same simultaneous coverage.
\section{An unbiased endpoint estimator without scalar exposure}
Allow arbitrary interference again. Suppose market assignments are independent across markets and endpoint probabilities $p_{m1},p_{m0}$ are positive and known. Under consistency,
\[
 \widehat\tau=\sum_m\nu_m\left[
 \frac{\ind\{Z_m=\mathbf1\}Y_m(Z_m)}{p_{m1}}
 -\frac{\ind\{Z_m=\mathbf0\}Y_m(Z_m)}{p_{m0}}
 \right]
 \tag{4}
\]
is unbiased for (1). To prove this, condition on the full potential-outcome table. On an endpoint assignment its observed outcome equals that endpoint potential. Each indicator divided by its assignment probability has expectation one. Taking expectations term by term gives (1), independently of all intermediate outcomes.

The endpoint events are mutually exclusive, so a single market's unweighted contribution has variance
\[
 \frac{Y_m(\mathbf1)^2}{p_{m1}}+
 \frac{Y_m(\mathbf0)^2}{p_{m0}}
 -\{Y_m(\mathbf1)-Y_m(\mathbf0)\}^2.
 \tag{5}
\]
Independence across markets makes total variance the sum of (5) times $\nu_m^2$. The unobserved joint endpoint pair means its exact numerical value is generally unavailable from a single realization. It is a valid design-variance formula, not a feasible variance estimator without further work.

For $n$ sellers independently treated with probability $1/2$, both endpoint probabilities equal $2^{-n}$. If both endpoint potential outcomes are one, the market contribution has mean zero and variance $2^{n+1}$. This gives exponential variance growth despite positive support. Randomizing whole markets directly to all-treated or all-control gives endpoint probabilities $1/2$ and variance $4$ in the same example. Increasing the number of seller observations inside a market is therefore not interchangeable with adding independent endpoint-assigned markets.

One can also see the practical rarity directly: with $M_0$ independent equal-sized markets, the probability of observing no all-treated market is $(1-2^{-n})^{M_0}$. When $M_0$ is much smaller than $2^n$, an unbiased estimator often receives no treated-endpoint contribution. Unbiasedness alone is a poor summary of its usefulness.
\section{What the proof does not identify}
The scalar exposure mapping is not established by randomization. Treating a large seller can alter the market differently from treating a small seller even when the fractions match. Equilibrium switching can violate a proposed Lipschitz constant, and outcome horizons can change the response curve. A one-month curve does not automatically describe welfare after years of entry.

Total transactions or revenue are also not total surplus. Demand, real costs, and the ownership boundary must justify the outcome measure. Consumer benefits may include participants who leave the platform, and ignoring them changes the population.

The exact interpolation result and randomization calculation are standard techniques supplied in full. The broad catalogue question remains unresolved because the relevant welfare measurement, exposure restrictions, equilibrium selection, and useful number of market experiments have not been established for an actual platform.
\begin{thebibliography}{9}
\bibitem{f} C. Farronato (2025). \emph{Digital Platforms as Information Aggregators: Implications for their Design and Regulation}, conference draft, section 1.1.
The primary source motivates interference and equilibrium adjustment; it does not supply the response restrictions imposed here.
\url{https://conference.nber.org/confer/2025/DTs25/farronato.pdf}.
\end{thebibliography}

\end{document}
